\begin{abstract}
Accurate mortality risk stratification from electronic health records demands
models that reason over longitudinal disease trajectories, numeric biomarkers,
and demographic context.
The central design question is \emph{how temporal structure should be encoded}:
does explicit geometric separation of events in time improve survival
discrimination over prose serialisation?
We present a controlled study across four input-representation families
(five experimental configurations) on a UK Biobank respiratory disease
subgroup ($n\,{=}\,124{,}158$), systematically varying input representation
from cross-sectional bag-of-features to unified clinical prose, with encoder
and fusion ablations across 12 experimental cells.
Our controlled design isolates input serialisation as the primary source of
variation: decoupled age--event sinusoidal encoding (Setting~4) achieves a
small but statistically robust $+$0.008 C-index advantage over prose
serialisation of identical events (bootstrap 95\,\% CI: $+$0.005 to
$+$0.010, $p<0.001$), while a Qwen3-8B prose embedding with no raw features
matches a 39-feature CoxPH baseline, demonstrating implicit numeric
precision in LLM representations.
Encoder choice and fusion method are secondary: within each setting, encoder
family changes C-index by at most 0.007 (and by only 0.001 in Setting~4),
while concatenation outperforms summation by $+$0.028--$+$0.039.
In zero-shot deployment, autoregressive generation shows a calibration gap
(IBS\,=\,0.185 vs.\ 0.141; AUC\,=\,0.52 at the earliest evaluated horizon,
9.5\,yr); this gap partly reflects missing acute biomarker inputs rather than
a fundamental limitation of the autoregressive paradigm.
\keywords{survival analysis \and LLM embeddings \and EHR \and
          UK Biobank \and temporal encoding}
\end{abstract}
